\chapter{Additional Derivation Details}
\label{app:derivation-details}

This appendix records algebra that supports Chapters~\ref{chap:baselines}
and~\ref{chap:expanded} but is not required to follow the main argument.

\section{Gradient of mutual information on the probability simplex}

Mutual information can be written as
\begin{equation}
  I(P)
  =\sum_{i,j}p_{ij}\log p_{ij}
  -\sum_i p_{i+}\log p_{i+}
  -\sum_j p_{+j}\log p_{+j}.
  \label{eq:appendix-mi-entropy-form}
\end{equation}
Treating the cell probabilities as coordinates gives
\begin{align}
  \frac{\partial I(P)}{\partial p_{ij}}
  &=(\log p_{ij}+1)
    -(\log p_{i+}+1)
    -(\log p_{+j}+1)\\
  &=\log\frac{p_{ij}}{p_{i+}p_{+j}}-1\\
  &=\pmi_P(i,j)-1.
  \label{eq:appendix-mi-gradient}
\end{align}
Any valid probability perturbation $h=(h_{ij})$ satisfies
$\sum_{i,j}h_{ij}=0$. Its directional derivative is therefore
\begin{align}
  \left.\frac{\mathrm d}{\mathrm d\varepsilon}
  I(P+\varepsilon h)\right|_{\varepsilon=0}
  &=\sum_{i,j}\{\pmi_P(i,j)-1\}h_{ij}\\
  &=\sum_{i,j}\pmi_P(i,j)h_{ij}.
  \label{eq:appendix-mi-directional}
\end{align}
The constant in the unconstrained gradient disappears on the simplex. For
the one-cell direction $h_{ij}=\mathbf 1\{i=x,j=y\}-p_{ij}$, Equation
\eqref{eq:appendix-mi-directional} becomes
\begin{equation}
  \pmi_P(x,y)-\sum_{i,j}p_{ij}\pmi_P(i,j)
  =\pmi_P(x,y)-I(P),
\end{equation}
which is Equation~\eqref{eq:mi-derivative}.

\section{Why the variance of pointwise MI appears}

For empirical cell probabilities,
\begin{equation}
  \widehat p_{ij}-p_{ij}
  =\frac{1}{n_P}\sum_{a=1}^{n_P}
  \{\mathbf 1(Z_a^{(P)}=(i,j))-p_{ij}\}.
\end{equation}
Substitution into the first-order expansion of MI gives
\begin{align}
  \Ihat(P)-I(P)
  &\approx
  \sum_{i,j}\pmi_P(i,j)(\widehat p_{ij}-p_{ij})\\
  &=\frac{1}{n_P}\sum_{a=1}^{n_P}
  \left[
  \sum_{i,j}\pmi_P(i,j)
  \{\mathbf 1(Z_a^{(P)}=(i,j))-p_{ij}\}
  \right]\\
  &=\frac{1}{n_P}\sum_{a=1}^{n_P}
  \{\pmi_P(Z_a^{(P)})-I(P)\}.
  \label{eq:appendix-mi-linear}
\end{align}
The observation-level summand has variance
\begin{align}
  \Var_P\{\pmi_P(X,Y)-I(P)\}
  &=\E_P\left[\{\pmi_P(X,Y)-I(P)\}^2\right]\\
  &=\sum_{i,j}p_{ij}\{\pmi_P(i,j)-I(P)\}^2\\
  &=V(P).
\end{align}
This is why the first-order variance of the complete MI estimator is
$V(P)/n_P$, rather than a variance of the numerical category labels.

\section{Mean of the variance influence function}

The influence function $g_P(x,y)$ is centred at the population level. This
can be shown without expanding its individual terms. The weighted average of
the one-cell perturbation directions is
\begin{equation}
  \sum_{x,y}p_{xy}\{\delta_{(x,y)}-P\}=P-P=0.
\end{equation}
Directional differentiation is linear in its direction. Hence
\begin{align}
  \E_P\{g_P(X,Y)\}
  &=\sum_{x,y}p_{xy}
    \left.\frac{\mathrm d}{\mathrm d\varepsilon}
    V\{(1-\varepsilon)P+\varepsilon\delta_{(x,y)}\}
    \right|_{0}\\
  &=D V_P\left[
    \sum_{x,y}p_{xy}\{\delta_{(x,y)}-P\}
    \right]\\
  &=D V_P[0]=0.
\end{align}
The empirical implementation nevertheless subtracts
$\overline g_P$ in Equation~\eqref{eq:tau-hat}. This removes floating-point
and plug-in centring error without changing the population definition.

\section{Direct table substitutions}

For observed counts $N^P=(N^P_{ij})$, the quantities used by expanded
Welch are calculated as
\begin{align}
  \widehat p_{ij}&=\frac{N^P_{ij}}{n_P},
  &\widehat p_{i+}&=\sum_j\widehat p_{ij},
  &\widehat p_{+j}&=\sum_i\widehat p_{ij},\\
  \widehat\pmi_P(i,j)
  &=\log\frac{\widehat p_{ij}}
  {\widehat p_{i+}\widehat p_{+j}},
  &\Ihat(P)&=\sum_{i,j}\widehat p_{ij}\widehat\pmi_P(i,j),\\
  \Vhat(P)&=\sum_{i,j}\widehat p_{ij}
  \{\widehat\pmi_P(i,j)-\Ihat(P)\}^2.
\end{align}
All sums involving pointwise MI are restricted to cells with
$\widehat p_{ij}>0$. The probability-weighted row and column means are
\begin{align}
  \widehat R_P(i)
  &=\frac{\sum_j\widehat p_{ij}\widehat\pmi_P(i,j)}
          {\widehat p_{i+}},\\
  \widehat C_P(j)
  &=\frac{\sum_i\widehat p_{ij}\widehat\pmi_P(i,j)}
          {\widehat p_{+j}}.
\end{align}
Consequently,
\begin{align}
  \widehat g_P(i,j)
  ={}&\{\widehat\pmi_P(i,j)-\Ihat(P)\}^2-\Vhat(P)\\
  &+2\{\widehat\pmi_P(i,j)-\widehat R_P(i)
  -\widehat C_P(j)+\Ihat(P)\},\\
  \overline g_P
  ={}&\sum_{i,j}\widehat p_{ij}\widehat g_P(i,j),\\
  \widehat\tau^2(P)
  ={}&\sum_{i,j}\widehat p_{ij}
  \{\widehat g_P(i,j)-\overline g_P\}^2,\\
  \widehat\nu_P
  ={}&\frac{2n_P\Vhat(P)^2}{\widehat\tau^2(P)}.
\end{align}
These expressions are the direct $O(rc)$ implementation of the population
derivation.

\section{Numerical derivative checks}

The implementation tests the analytic derivatives against centred finite
differences on positive probability tables. For a selected cell, define
\begin{equation}
  P_{\pm h}=(1\mp h)P\pm h\delta_{(x,y)}.
\end{equation}
For sufficiently small $h$, the checks compare
\begin{align}
  \frac{I(P_h)-I(P_{-h})}{2h}
  &\quad\text{with}\quad \pmi_P(x,y)-I(P),\\
  \frac{V(P_h)-V(P_{-h})}{2h}
  &\quad\text{with}\quad g_P(x,y).
\end{align}
The finite-difference step is selected so that both perturbed tables retain
positive probabilities. These checks test the algebra independently of the
Monte Carlo calibration experiment.
